\documentclass[10pt,twocolumn,a4paper]{article}

\usepackage[margin=0.85in,columnsep=0.28in]{geometry}
\usepackage[numbers]{natbib}
\usepackage{booktabs}
\usepackage{amsmath}
\usepackage{amssymb}
\usepackage{microtype}
\usepackage{url}
\usepackage[colorlinks=true,allcolors=blue]{hyperref}

\newcommand{\ebar}{\bar{\epsilon}}

\title{Measuring Agents by Their Trajectories:\\
Process Supervision and Its Verifier Dependency}

\author{Literature review}
\date{August 2026}

\begin{document}
\maketitle

\begin{abstract}
Language-model agents have become markedly better at completing multi-step tasks
without becoming better at establishing whether they completed them.
The dominant response across 2026 is to replace a single terminal reward with dense,
step-level signal, in both training and evaluation.
We review that turn across eighteen 2026 papers, anchored on three earlier works that set
the agenda, and argue three things.
First, the case against outcome-only measurement is no longer an inference but a quantity:
re-scoring three computer-use benchmarks with a purpose-built rubric verifier moves reported
success from $74.6\%$ to $37.9\%$ on one, with comparable gaps on every benchmark--agent pair
tested \citep{rosset2026artbuildingverifierscomputer}.
Second, dense supervision is conditional rather than free: a verifier with high
disagreement is measurably worse than no process signal at all, in-domain and
out-of-domain \citep{yuan2026verifiableprocessrewardsagentic}, and in multi-agent settings
process verification does not reliably help
\citep{venkataramani2026masproveunderstandingprocessverification}.
Third, and cutting across the rest, agent interventions move the quantity being measured
more reliably than the quantity being wanted --- a dense tool-use reward raises tool
adoption from $0.03$ to $0.33$ with no held-out accuracy gain
\citep{fan2026screenshotstoolselicitingtool}, self-evolution raises utility and
unauthorized state changes together \citep{li2026auditingselfevolutionfinancialagents}, and
self-authored skills score $8$--$11$ points \emph{below} using no skill
\citep{peng2026writeexecuterefineskill}.
We close with four scoped experiments that this corpus states or implies are missing, the
most consequential being to measure verifier disagreement on a realistic agentic task and
locate the threshold at which process supervision becomes harmful.
\end{abstract}

\section{Introduction}

A frontier language model now solves in a single forward pass problems that were recent
research contributions.
Embedded in an agent loop and asked to carry a task over many steps, the same model
fails in ways no single-step benchmark reveals.
\citet{chen2026horizongapplanningmemory} name this the \emph{horizon gap}: ``the distance
between what a model can do in one step and what a system built around it can reliably
finish over many'', with characteristic failures being ``losing track of an earlier
decision, declaring a half-finished job done, or quietly drifting from the goal they were
given.''

Two of those three failures are failures of self-assessment rather than of capability.
An agent that declares a half-finished job done has not run out of ability; it has run out
of the means to check.
This review takes that observation as its organizing thread, because it is also the shape
of the field's dominant 2026 research programme.

The programme has a common diagnosis, stated in near-identical language across
otherwise unrelated papers.
Outcome-only credit assignment gives ``the same credit to effective elicitation, errors,
and later repair'' \citep{zhao2026betteragentsmultiturnuser}; a trajectory ``may fail
despite containing many correct intermediate decisions, or succeed despite containing
flawed ones'' \citep{yuan2026verifiableprocessrewardsagentic}; most benchmarks ``judge only
final outcomes, unable to distinguish reliable reasoning from lucky success''
\citep{wu2026clawtracktracelevelevaluationimprovement}.
The proposed cure is uniformly the same: manufacture denser signal at the level of the
step or the trajectory segment.

The interesting finding of the past year is not that the cure works.
It is that the cure is \emph{conditional}, and that its precondition --- a verifier good
enough to be worth listening to --- is the quantity nobody measures.

\paragraph{Contributions.}
This review makes four.
(i) We separate three properties the literature routinely conflates, following
\citet{chen2026horizongapplanningmemory}, and use the distinction as the spine of the
review (\S\ref{sec:terms}).
(ii) We consolidate the evidence that outcome-only measurement systematically overstates
agent success, and characterise the four distinct mechanisms by which it does so
(\S\ref{sec:outcome}).
(iii) We organise the process-supervision literature by \emph{where its dense signal comes
from} rather than by method family, which makes the shared verifier dependency visible
(\S\S\ref{sec:process}--\ref{sec:condition}).
(iv) We identify a cross-cutting pathology that no single paper in the corpus names ---
interventions reliably move the measured quantity and unreliably move the intended one ---
and derive a research agenda from what the corpus states is missing
(\S\S\ref{sec:pathology}--\ref{sec:agenda}).

\section{Scope and Method}
\label{sec:scope}

This is a focused review, not a systematic one, and the selection procedure matters for
how its conclusions should be read.

\paragraph{Corpus.}
We reviewed eighteen papers posted to arXiv in 2026, selected by searching eight query
threads spanning process rewards, trajectory verification, mid-trajectory recovery,
long-horizon memory, agentic reinforcement learning, computer-use agents, deep-research
agents, and reliability measurement, then filtering to 2025-or-later and ranking by
relevance to five weaknesses already identified in earlier reading.
Three earlier works anchor the discussion because they set the terms the 2026 papers
respond to: ReAct \citep{yao2023reactsynergizingreasoningacting}, Reflexion
\citep{shinn2023reflexionlanguageagentsverbal}, and $\tau$-bench
\citep{yao2024taubenchbenchmarktoolagentuserinteraction}.

\paragraph{Reading depth.}
Seven of the eighteen were read in full; eleven at abstract depth, used only for claims
their abstracts state directly and never for numbers requiring a table, and flagged as such
at first substantive use.

\paragraph{Secondary citations.}
\citet{chen2026horizongapplanningmemory} is a survey of 1{,}547 papers, and several
striking figures in \S\ref{sec:outcome} and \S\ref{sec:repair} reach us through it rather
than from the primary source: the SWE-bench audits, the Vending-Bench coherence results,
the AndroidControl noise ceiling, and the accuracy--correction paradox.
We attribute these to the survey and do not treat them as independently verified; a reader
building on any should read the primary paper first.
A review that launders secondary citations into apparent primary evidence commits a smaller
version of the error this review is about.

\paragraph{What we do not cover.}
Planning, orchestration topologies, GUI grounding, and agent safety are touched only where
they bear on measurement; \citet{chen2026horizongapplanningmemory} covers the full pipeline
at survey scale.

\section{Three Properties, Routinely Conflated}
\label{sec:terms}

\citet{chen2026horizongapplanningmemory} open by fixing vocabulary, and the distinction is
load-bearing enough to restate.
Three properties are logically independent:

\begin{itemize}
\item \textbf{Long-context} is a property of the \emph{model}: how many tokens one forward
      pass can attend to.
\item \textbf{Long-horizon} is a property of the \emph{task}: how many sequential decisions
      it requires, regardless of whether the resulting trace fits in a context window.
\item \textbf{Long-term memory} is a property of the \emph{system}: whether information
      available at step $t$ remains available at step $t+k$.
\end{itemize}

The independence is not pedantic: ``a harness can carry a long-horizon task through a
short context via aggressive summarization, and a system can have a huge context window
and no persistent memory at all''
\citep{chen2026horizongapplanningmemory}.
The cost of conflating them is that ``a `memory' paper and a `long-context' paper are so
often cited as if they addressed the same problem when they do not.''

Two further distinctions recur below.
The \textbf{agent} is the policy --- the model plus whatever turns its output into actions
--- while the \textbf{harness} is the surrounding software: control loop, tool sandbox,
memory layer, retry logic, sub-agent orchestration.
An \textbf{episode} is one complete attempt at a task; a \textbf{session} is a continuous
interactive period bounded by the interface rather than the task, so a session may contain
many episodes and an episode may span sessions.

The agent/harness split is where this review's most uncomfortable open question lives
(\S\ref{sec:agenda}).
Relatedly, the locus at which extra horizon is carried determines the failure mode:
within-context methods fail by attention degradation, methods exceeding the window inside
one episode fail by ``information loss at the harness boundary'', and methods persisting
across episodes fail by ``interference or staleness in what was retained''
\citep{chen2026horizongapplanningmemory} --- so treating memory as one bucket obscures three
unrelated mechanisms.

\section{What Outcome-Only Measurement Hides}
\label{sec:outcome}

\subsection{The admission that started it}

\citet{yao2024taubenchbenchmarktoolagentuserinteraction} introduced $\mathrm{pass}^k$, ``the
chance that all $k$ i.i.d. task trials are successful'', reporting ``$<50\%$'' at
$\mathrm{pass}^1$ and ``$\mathrm{pass}^8 < 25\%$'' on retail.
More consequentially, they stated the limitation of their own terminal-state reward:
``$r = 1$ might be a necessary but not sufficient condition for a successful episode e.g.,
the agent might issue the return without explicit user confirmation, which violates the
policy.''
A trajectory can violate the domain policy and still score $1$.

That was an inference about one benchmark.
It is now a measured, cross-benchmark quantity.

\subsection{The size of the gap}

\citet{rosset2026artbuildingverifierscomputer} built a rubric-based verifier for
computer-use trajectories --- approximately $3{,}000$ lines of code and $2{,}000$ lines of
prompts, hand-iterated --- and re-scored three benchmarks with it.
Table~\ref{tab:gap} reports the result.

\begin{table}[t]
\centering
\begin{tabular}{@{}llrr@{}}
\toprule
Benchmark & Agent & Native & Rubric \\
\midrule
WebVoyager        & Fara-7B & $74.6$ & $37.9$ \\
WebVoyager        & GPT-5   & $90.6$ & $71.0$ \\
Online-Mind2Web   & Fara-7B & $32.2$ & $15.8$ \\
Online-Mind2Web   & GPT-5   & $62.0$ & $48.6$ \\
WebTailBench      & Fara-7B & $39.6$ & $23.2$ \\
WebTailBench      & GPT-5   & $62.5$ & $39.9$ \\
\bottomrule
\end{tabular}
\caption{Reported success rate (\%) under each benchmark's own (\emph{native}) verifier
versus the \emph{rubric} verifier's outcome judgement, from
\citet{rosset2026artbuildingverifierscomputer}.
False positive rates of native verifiers against the rubric verifier exceed $20\%$ on
every pair.
The rubric verifier is treated as the reference label, so these are disagreements rather
than proven errors --- see the caveat in \S\ref{sec:calib}.}
\label{tab:gap}
\end{table}

Reaching those numbers required solving the verification problem well enough to be
credible, and the paper is candid that this is hard.
Its best outcome agreement with human labels is Cohen's $\kappa$ of $0.64$ and $0.58$ on
two datasets, against baselines at $0.44$/$0.26$ and $0.31$/$0.13$, and against human
inter-annotator agreement of $0.53$--$0.57$.
The prior bound it cites is that ``no LLM-based judge exceeds $70\%$ precision'' against
human agreement of $89.3\%$.
Best-in-class verification is therefore roughly as good as humans who agree with each
other poorly.

\subsection{Four distinct mechanisms}

The corpus supports separating four ways outcome-only measurement misleads, which are
usually run together under ``benchmarks are flawed''.

\paragraph{(1) Leniency.}
The classic case: the check passes on trajectories that should fail.
As reported by \citet{chen2026horizongapplanningmemory}, an audit of SWE-bench found
$32.67\%$ of successful patches involved solution leakage and a further $31.08\%$ passed only
because the test suite was too weak to verify correctness --- filtering both dropped one
leaderboard-topping system from $12.47\%$ to $3.97\%$.
Across agentic benchmarks generally, setup and reward-design flaws ``can over- or
under-estimate reported performance by up to'' $100\%$ ``in relative terms''.

\paragraph{(2) Noise as a ceiling.}
The inverse failure, and less discussed.
Top-scoring GUI agents ``plateau around'' $60\%$ on AndroidControl ``because benchmark
noise --- ambiguities and factual errors in the benchmark itself --- caps the achievable
score'' \citep{chen2026horizongapplanningmemory}.
This has a three-year-old precedent: \citet{yao2023reactsynergizingreasoningacting} report
label ambiguity at $29\%$ of judged cases against an exact-match gap to chain-of-thought of
$27.4$ versus $29.4$, placing the headline comparison inside annotation noise.

\paragraph{(3) Harness artifacts.}
The most alarming instance we found is not about agents at all.
Auditing three self-evolution methods,
\citet{li2026auditingselfevolutionfinancialagents} found that one carried a literal
WebArena text-action envelope that disrupted tool execution in their function-calling
executor.
Removing only that envelope moved utility from $0.319$ to $0.756$.
A $43$-point swing from a prompt-format mismatch is larger than any effect that study was
designed to measure, and it means cross-harness comparisons of these methods may be
measuring envelope compatibility.
The paper's conclusion --- that auditing requires tracking ``artifact-executor
compatibility'' --- generalises beyond its domain.

\paragraph{(4) Whether the run finished.}
\citet{khanal2026pass1reliabilityscienceframework} argue that completion rate should be
``a first-class validity metric'', having lost $368$ episodes between $23{,}760$ planned and
$23{,}392$ completed ``after deduplication'' without stating what was removed.
\citet{rosset2026artbuildingverifierscomputer} report $17.0\%$ unterminated trajectories
for one agent on WebTailBench and $7.2\%$ for another on Online-Mind2Web, with no statement
of how they count toward reported success.
A $17\%$ unaccounted slice is larger than several of the disagreements being interpreted in
the same table.

\subsection{The gap is real but uncalibrated}
\label{sec:calib}

Two facts prevent Table~\ref{tab:gap} from being an estimate of overstatement.
The rubric verifier is the reference label with its own human agreement at
$\kappa \approx 0.6$; and its near-zero false positive rate ($0.01$, $0.08$) is an
optimisation constraint rather than a discovery --- the objective was to ``maximize Cohen's
$\kappa$ without increasing FPR; any FPR-increasing change is automatically rolled back'' ---
with a corresponding false \emph{negative} rate of $0.31$--$0.32$.
A verifier tuned never to over-credit under-reports success, the same direction that makes
native verifiers look loose.
Benchmark verifiers are \emph{loose} by an amount this literature has not calibrated.

\section{The Turn to Process Supervision}
\label{sec:process}

\subsection{Organising by signal source, not method family}

Process-reward papers are conventionally grouped by algorithm.
Grouping them by \emph{where the dense signal originates} is more informative, because it
exposes what each one needs from its environment.
Table~\ref{tab:sources} does this for the corpus.

\begin{table*}[t]
\centering
\begin{tabular}{@{}p{4.1cm}p{3.6cm}p{5.6cm}l@{}}
\toprule
Signal source & Instance & Requires & Depth \\
\midrule
Symbolic/algorithmic oracle & VPR \citep{yuan2026verifiableprocessrewardsagentic}
  & a solver or search per environment & full \\
Human step annotation & AgentProcessBench \citep{fan2026agentprocessbenchdiagnosingsteplevelprocess}
  & per-step human labels & full \\
Rubric + judge over screenshots & Universal Verifier \citep{rosset2026artbuildingverifierscomputer}
  & rubric generation, judge model & full \\
Rubric over reasoning turns & ClawTrack \citep{wu2026clawtracktracelevelevaluationimprovement}
  & task-specific rubric items & abstract \\
Environment's next reaction & Faca \citep{zhao2026betteragentsmultiturnuser}
  & an instrumented simulator & full \\
Privileged self-teacher & CrEST \citep{wang2026teachmagnitudedirectionverifierbounded}
  & a stronger view of the same policy & abstract \\
Policy's own belief consistency & ReBel \citep{tang2026rewardingbeliefsactionsconsistencyguided}
  & nothing external & abstract \\
\bottomrule
\end{tabular}
\caption{Sources of dense step-level signal in the reviewed corpus, ordered by decreasing
external requirement.
``Depth'' records whether we read the paper in full or at abstract depth
(\S\ref{sec:scope}).
The trend down the table is the story: each row buys density at lower external cost, and
no paper establishes where signal quality falls off.}
\label{tab:sources}
\end{table*}

\subsection{The strongest positive result}

\citet{zhao2026betteragentsmultiturnuser} give the cleanest demonstration that dense credit
beats terminal credit, and the reason is experimental design rather than effect size.
Their control shares the frozen simulator, the agent-visible dialogue, the supervised
initialisation, the training data, the rollout construction, the optimiser, and the
horizon; credit assignment is the only difference.
Under that control the nine-domain $\tau$-family average moves from
$34.66 \pm 0.25$ to $40.57 \pm 1.04$ at 8B and from $42.51 \pm 0.12$ to
$52.73 \pm 1.53$ at 14B, with the ordering holding across three training seeds at both
scales.

Two qualifications matter.
The gains concentrate sharply --- at 14B, two of nine cells move from $26.3$ to $83.6$ and
from $50.6$ to $82.7$ --- and the signal is the simulator's \emph{private} strategy label,
which real users do not emit, so the deployable variant is explicitly future work.
The polarity map moreover ``describe[s] interaction movement rather than user sentiment or
action correctness'', so a policy-violating action a satisfied user confirms receives
positive credit.
$\tau$-bench's complaint was that terminal reward cannot see policy compliance; this adds
dense credit that also cannot.

\subsection{The label schema is where the gaps are}

\citet{fan2026agentprocessbenchdiagnosingsteplevelprocess} contribute the first
human-annotated step-level benchmark for tool-using agents: $1{,}000$ trajectories over
$200$ unique tasks, $8{,}509$ annotated actions, inter-annotator agreement $89.1\%$ and
Cohen's $\kappa$ of $0.767$, with ternary labels and an error-propagation rule marking
everything causally downstream of a mistake as incorrect.

Three structural observations follow, in rising order of consequence.

First, the neutral label --- introduced precisely so exploratory steps are not punished ---
is a small share of steps and the hardest class to predict, so under a micro-averaged
step-accuracy metric a model that never predicts it forfeits almost nothing.
The headline ranking is close to blind to the property the benchmark was built to measure.

Second, the annotation protocol moves results more than the model does.
Removing error propagation changes downstream Best-of-8 accuracy by $+3.7$ for one
evaluator and by $-13.2$ for a \emph{variant of the same model} ($64.2$ to $50.9$).
Anything built on this benchmark inherits a protocol choice whose effect size rivals the
thing being measured.

Third, and most attackable: irreversibility motivates the benchmark and appears nowhere in
it.
The paper opens on the claim that ``tool execution frequently entails irreversible side
effects'' --- the stated reason step-level verification matters more for agents than for
mathematics --- yet no label marks reversibility.
A formatting error and an irreversible deletion are both scored $-1$.
A process reward model trained on these labels learns \emph{wrong}, never
\emph{unrecoverable}, which is the distinction an agent would need in order to stop and ask.

\section{The Reliability Condition}
\label{sec:condition}

\subsection{A bound and an ablation}

\citet{yuan2026verifiableprocessrewardsagentic} give the clearest treatment because they
can drive verifier error to zero by construction, using symbolic oracles in Tic-Tac-Toe,
Sudoku, and Minesweeper.
Their analysis bounds the policy-gradient bias by the verifier disagreement rate:
$\lVert \hat{g} - g^{\star} \rVert \le G\ebar$, so ``oracle error propagates one-to-one
into the gradient, with no horizon-dependent amplification.''

The ablation matters more than the headline.
Varying the number of Monte Carlo tree search simulations in the oracle yields
Table~\ref{tab:oracle}.

\begin{table*}[t]
\centering
\begin{tabular}{lcc}
\toprule
Oracle strength & In-domain return (1st / 2nd) & OOD average \\
\midrule
Base model (no training) & $-0.31 \pm 0.04$ / $-0.35 \pm 0.05$ & $60.92$ \\
$N = 100$                & $-0.48 \pm 0.06$ / $-0.52 \pm 0.07$ & $58.47$ \\
$N = 1000$               & $-0.13 \pm 0.04$ / $-0.15 \pm 0.04$ & $61.86$ \\
$N = 10{,}000$ (default) & $-0.09 \pm 0.03$ / $-0.11 \pm 0.03$ & $62.17$ \\
\bottomrule
\end{tabular}
\caption{Sensitivity of verifiable process rewards to oracle quality
\citep{yuan2026verifiableprocessrewardsagentic}.
In-domain return is on Tic-Tac-Toe (optimum $0$); the out-of-domain column is a
seven-benchmark general-reasoning average.
A weak oracle is worse than not training at all, on both axes.}
\label{tab:oracle}
\end{table*}

The $N = 100$ row is the load-bearing result of this entire literature.
Dense supervision from a weak verifier is not merely ineffective; it is net-harmful
relative to leaving the pretrained policy alone, with degradation across every downstream
benchmark.
The authors' own summary: ``dense feedback alone is not sufficient: for process rewards to
improve long-horizon reasoning, they must also be reliable and objectively grounded.''

Two caveats bear on how far the bound has been tested: the default oracle runs $10{,}000$
simulations per move against the rollout baseline's $100$ rollouts per state, so the
comparison is not budget-matched; and the Sudoku verifier reads the unique solution grid,
making that instantiation filtered imitation on labelled solutions rather than verification
of task structure.

\subsection{And in multi-agent settings it does not work}

\citet{venkataramani2026masproveunderstandingprocessverification}, read at abstract depth,
run the comparison at scale --- three verification paradigms, two granularities, five
verifiers, four context-management strategies, six multi-agent frameworks --- and report
that ``process-level verification does not consistently improve performance and frequently
exhibits high variance, highlighting the difficulty of reliably evaluating partial
multi-agent trajectories.''

Their most consequential sentence receives one clause: ``a small performance gap between
LLMs acting as judges and as single agents.''
If a model is no better at judging a trajectory than at producing one, process verification
reduces to an ensemble technique, and the premise beneath this research programme is weaker
than it appears.

\subsection{Synthesis}

Taken together the picture is coherent and uncomfortable.
Process supervision helps when the verifier is good, is harmful when the verifier is bad,
has a theoretical account of the transition, and has no measurement of where any realistic
verifier sits.
The best hand-built verifier for a realistic domain reaches $\kappa \approx 0.6$
\citep{rosset2026artbuildingverifierscomputer}; the harmful regime is demonstrated at
$N=100$ \citep{yuan2026verifiableprocessrewardsagentic}; nothing connects the two scales.
``Add process supervision'' is therefore a bet, not an improvement.

\section{Repair Within the Trajectory}
\label{sec:repair}

\subsection{From restart to repair}

For two years the literature offered exactly two responses to a derailed trajectory.
ReAct exhibits the first --- continue and fail, with loop entrapment where ``the model
repetitively generates the previous thoughts and actions'' and non-informative search at
$23\%$ of error cases, made terminal by a hard step budget of $7$ steps on HotpotQA and $5$
on FEVER \citep{yao2023reactsynergizingreasoningacting}.
Reflexion institutionalised the second: abandon the episode and restart from a reset
environment \citep{shinn2023reflexionlanguageagentsverbal}.

\citet{yang2026selfcorrectinglonghorizonsearchagents} supply the missing third option.
Their memory is a dependency tree in which ``a child is a search state derived from its
ancestor's evidence'' --- explicitly not a search over candidate continuations.
On a confirmed contradiction the system locates the node that introduced the refuted fact,
replaces the fact and its source, regenerates that node's summary from corrected evidence,
prunes all descendants, and resumes.
This instantiates dependency-directed revision from truth-maintenance systems inside an
agent loop.

Their framing identifies precisely what verbal reflection cannot do: existing compression
methods often ``replace erroneous facts without repairing downstream reasoning derived from
them.''
Fixing the belief is not fixing the trajectory.

Across $2{,}149$ questions on four search benchmarks, judge accuracy improves from $30.1$
for full-trajectory ReAct to $44.0$, a pooled gain of $13.9$ points and $8.3$--$25.6$ points
per dataset.

\subsection{Why that result is not yet attributable}

The comparison isolating the repair mechanism is against a flat-memory control sharing the
same summary budget, evidence budget, and conflict detector: $2.2$--$4.7$ points.
But backtracking ``is triggered in $9.6$--$17.5\%$ of ReTree runs'', and on the rest the two
systems should behave near-identically --- so the whole gap must sit in the minority that
backtracked, requiring the repair to flip a fifth to a third of them.
The paper never conditions accuracy on whether a backtrack fired, and uses a single seed with
no confidence intervals.
Stratifying already-collected data by that condition is the cheapest informative experiment
we identify.

Two further limits are worth recording.
Peak per-step context for the full-trajectory baseline is $1{,}920$ characters, far below
any context limit, so whatever the method fixes is not a context-capacity effect despite
being motivated as one.
And the provenance metric the architecture exists to deliver is citation precision of
$44.5$: most cited passages do not entail the claim attached to them.

\subsection{The compounding model that motivates repair is wrong}

Repair mechanisms are designed against an intuition of independent per-step error
compounding exponentially in horizon.
The measurements do not support it.

As reported by \citet{chen2026horizongapplanningmemory}, Vending-Bench --- on runs
exceeding $20$M tokens --- finds not smooth decay but that ``capable models turn a profit
in most runs, but every model has some runs that derail into a `meltdown' loop from which
they rarely recover'', and critically that ``derailment shows no clear correlation with the
context window filling up.''
Their synthesis is that degradation is bimodal, ``(fine, or catastrophically derailed)
whose trigger the field cannot yet reliably predict''.

\citet{khanal2026pass1reliabilityscienceframework} reach a compatible conclusion from a
different direction.
Across $396$ tasks, four duration buckets, ten models, and $23{,}392$ episodes, aggregate
$\mathrm{pass}@1$ falls from $76.3\%$ to $52.1\%$, a $24.3$-point decline that runs faster
than an i.i.d. Bernoulli baseline: for one model the geometric prediction is
$0.758^4 \approx 33.0\%$ against $22.2\%$ observed.
The mechanism is positive inter-step error correlation, with the theoretical consequence
that failure scales as $\Omega(\epsilon \cdot e^{\rho T})$, so ``Reducing $\rho$ \ldots is
therefore a higher-priority training objective than reducing $\epsilon$ alone''.

That is the strongest available argument for repair as a research direction, and it
reframes the target: the quantity to attack is error \emph{persistence}, not error rate.
It also implies that step budgets and context-compaction thresholds justified by
exponential-decay reasoning are fitted to a curve the data rejects.

\subsection{Self-correction may not improve with scale}

A counter-intuitive result deserves recording, with the caveat that it reaches us
second-hand and rests on three models.
Decomposing self-correction into detection, localisation, and correction yields an
``accuracy-correction paradox'': the weaker model ($66\%$ base accuracy) corrects its own
errors intrinsically at $26.8\%$ while the strongest ($94\%$ base accuracy) manages
$16.7\%$, and detection rate does not predict correction success
\citep{chen2026horizongapplanningmemory}.
The proposed Error Depth Hypothesis --- stronger models make fewer but structurally deeper
errors that intrinsic correction cannot reach --- implies the long-standing negative result
on unaided self-correction ``could get worse, not better, as models improve.''

Consistent with this, \citet{yu2026trainingfreeinferencetimeselfreflectioncostbounded},
at abstract depth, report a training-free generate--critique--revise loop that ``does not
raise accuracy beyond the $95\%$ Wilson interval'', repositioning its contribution as cost
control ($82$--$88\%$ of items stopped early at equal accuracy) --- and it is the cheap
baseline trained methods rarely compare against.

\section{Accumulation and Its Costs}
\label{sec:accum}

Accumulation operates across episodes where repair operates within one, on the premise that
experience compounds into reusable skills, workflows, or memories.
Three results indicate the premise is not free.

\paragraph{Accumulated experience can be worse than none.}
\citet{peng2026writeexecuterefineskill} open on the observation that ``agent-authored
skills perform $8$--$11$ points worse than using no skill'', against expert-written skills
that do help, concluding that following procedural guidance and improving it from execution
evidence are distinct capabilities.
Their remedy trains a skill optimiser outside a frozen executor with a programmatic verifier
scoring repeated executions, gaining $+7.80$ and $+3.85$ points over no-skill on two
benchmarks --- so on one it recovers a self-authoring deficit and adds under four points,
and it needs a verifier its motivating case of novel software lacks.
Architecturally it nonetheless answers an open question: Reflexion's retries occur at
deployment, here they occur during optimiser training, so the deployed executor is
single-pass --- the transfer from resettable environment to single-pass deployment that
$\tau$-bench's setting demands.

\paragraph{Capability and safety move on different axes.}
\citet{li2026auditingselfevolutionfinancialagents} audit three published self-evolution
methods in simulated e-banking, tracking regressions, attack-surface contact, unauthorized
state change, and artifact compatibility rather than accuracy alone.
For one method, benign utility rises from $0.741$ to $0.837$ while exposure to injected
content rises from $0.820$ to $0.943$; conditional attack success after exposure
\emph{falls} from $0.605$ to $0.562$ while overall attack success \emph{rises} from
$0.496$ to $0.530$, with unauthorized financial state changes at $0.685$.

The mechanism is exposure rather than vulnerability: the evolved agent is individually more
robust per encounter and more successfully attacked overall, because it reaches more
injected content.
A paper reporting only conditional robustness would have shown a safety improvement from
the same run.
The general form, in the authors' words, is that ``post-evolution accuracy alone does not
show whether learned behavior preserves previously correct behavior or security.''

\paragraph{Memory scaffolds do not help at long horizons.}
The adjacent intervention fails outright.
\citet{khanal2026pass1reliabilityscienceframework} compare a reason-act loop against the
same loop with an episodic scratchpad across ten models: six are hurt, four neutral within
$\pm 0.03$, largest penalties $-0.14$ and $-0.13$, with the effect concentrated in models
``capable enough to use the scratchpad but not capable enough to absorb its overhead
efficiently.''
This is in tension with the Reflexion lineage, where an episodic buffer is the mechanism
carrying learning across trials; the settings differ --- between-episode lessons versus
within-episode notes --- but the result should temper any default use of episodic memory as
a reliability intervention.

\section{A Cross-Cutting Pathology}
\label{sec:pathology}

No single paper in the corpus names the following, but four exhibit it, and we regard it as
the review's most portable observation.

\textbf{Agent interventions move the quantity being measured more reliably than the
quantity being wanted.}

\begin{itemize}
\item \citet{fan2026screenshotstoolselicitingtool}: under a fixed harness, a dense tool-use
      reward raises spreadsheet adoption from $0.03$ to $0.33$ and carries into greedy
      decoding, but ``held-out accuracy does not follow. Behavior is steerable; competence
      is not.''
      The reward is entirely benign, making this structural rather than a reward-design
      error.
\item \citet{li2026auditingselfevolutionfinancialagents}: utility and unauthorized state
      change rise together, and only the first is what the method was reported on.
\item \citet{peng2026writeexecuterefineskill}: self-authored skills look like
      self-improvement and score below no skill at all.
\item \citet{wu2026clawtracktracelevelevaluationimprovement}, at abstract depth: outcome
      scoring cannot separate reliable reasoning from ``lucky passes'', and identifies
      \emph{result verification} as the systematic bottleneck among the process dimensions
      it scores --- an agent that does not check its own results cannot detect the state a
      repair would act on.
\end{itemize}

The same paper that produces the sharpest negative result also produces the corpus's
cleanest positive one, and the contrast is instructive.
Incentivising what the agent \emph{does} failed; changing what it \emph{observes} worked ---
dropping the screenshot rendered redundant by a successful tool call and retraining under
that rule reaches $37.8\%$ against $33.0\%$ at $53\%$ of the input cost
\citep{fan2026screenshotstoolselicitingtool}.

The same paper also demonstrates that the harness-versus-model question has no
model-independent answer: under one identical harness on $309$ tasks, the same tool
availability improves a reasoning model by $+4.0$ points and degrades a non-reasoning model
by $-5.9$ points, both beyond two standard errors.
And the reasoning model that benefits calls a tool on only $55$ of $309$ tasks, $23.9\%$ of
the tool-reachable ones --- an adoption gap the authors attribute to the model having ``a
cheaper route'' and never being trained to take it.

\citet{chen2026horizongapplanningmemory} name the structural risk this implies but do not
measure it: the process signals used to \emph{train} long-horizon agents and those used to
\emph{evaluate} them ``rest on overlapping assumptions about what counts as progress on a
partial trajectory'', and the disjoint-signal comparison that would test it is ``an
experiment the corpus does not currently contain.''

\section{Research Agenda}
\label{sec:agenda}

We list four experiments the reviewed corpus states or implies are missing, ordered by
cost against expected value.
Each is smaller than a paper.

\paragraph{1. Measure verifier disagreement on a realistic agentic task, and locate the
harmful threshold.}
Both halves already exist:
\citet{rosset2026artbuildingverifierscomputer} supply a strong verifier and human-labelled
trajectories, and \citet{yuan2026verifiableprocessrewardsagentic} supply the bound
$\lVert \hat{g} - g^{\star} \rVert \le G\ebar$ together with a demonstration of the harmful
regime.
Nothing joins them.
Estimating $\ebar$ for obtainable verifiers, and finding where process supervision crosses
from helpful to harmful, would convert this literature's central assumption into a measured
quantity.

\paragraph{2. For any dense agent reward, check whether competence followed behaviour.}
\citet{fan2026screenshotstoolselicitingtool} is the only paper in the corpus that runs this
check, and it fails.
Re-running held-out evaluation across published process-reward methods is cheap and, on
this evidence, likely to find several that move the proxy without moving the capability.

\paragraph{3. Stratify repair results by whether repair fired.}
\citet{yang2026selfcorrectinglonghorizonsearchagents} report that backtracking triggers in
$9.6$--$17.5\%$ of runs and never separate those runs from the rest.
The split requires no new experiments and determines whether dependency-directed revision
does anything.

\paragraph{4. Encode reversibility.}
Irreversibility is the stated motivation for step-level verification and appears in no
label schema, reward, or metric we reviewed.
Adding a reversibility flag to the existing $8{,}509$ annotations of
\citet{fan2026agentprocessbenchdiagnosingsteplevelprocess} requires no new trajectories.
The stronger version --- a reward weighting errors by recoverability, and a benchmark that
credits an agent for declining an irreversible action under uncertainty --- is closer to a
thesis than a paper.

Prior to most method work sits one attribution question.
\citet{chen2026horizongapplanningmemory} call separating the model's own competence from the
harness compensating for its absence among ``the field's most consequential open measurement
problems'', with ``the corpus does not yet contain the controlled comparison that would answer
it.''
Since \citet{fan2026screenshotstoolselicitingtool} show the interaction term is significant,
every single-model scaffold result here is uninterpretable as a claim about scaffolds.

\section{Threats to Validity}
\label{sec:threats}

This review has the weaknesses of its construction, and they run in identifiable
directions.

\paragraph{Selection.}
Relevance-ranked search weighted toward five previously identified weaknesses biases the
corpus toward papers that diagnose problems over papers that quietly solve them, and
under-represents engineering that does not frame itself as critique.
Our conclusion that the field's dominant activity is measurement repair should be read with
that in mind.

\paragraph{Reading depth.}
Eleven of eighteen papers were read at abstract depth (Table~\ref{tab:sources} records
which).
The failure modes visible only on close reading --- a calibrated threshold that turns out to
be zero, a metric defined one way and computed another --- are exactly what abstract-depth
reading misses; each of the seven full reads surfaced at least one, which suggests the
eleven contain more.

\paragraph{Secondary citation and recency.}
Several of our most quotable figures pass through
\citet{chen2026horizongapplanningmemory} rather than the primary source
(\S\ref{sec:scope}); we mark but do not verify these.
The corpus is almost entirely unrefereed 2026 preprints, effect sizes in preprints tend to
shrink under review, and several results rest on a single seed or model.

\section{Conclusion}

The reviewed literature is best read as a field discovering that its instruments are the
bottleneck.
Agents improved at completing tasks faster than the means of establishing completion
improved, and the response --- denser signal at the step and segment level --- is being
pursued simultaneously in training and evaluation, from shared intuitions about what partial
progress looks like.

Three findings should survive.
Outcome-only measurement overstates agent success by a large, uncalibrated margin through at
least four distinct mechanisms.
Dense process supervision is a bet on verifier quality with a known harmful regime and no
measurement of where realistic verifiers fall.
And interventions move the measured quantity considerably more reliably than the intended
one, which makes the default protocol --- optimise a proxy, report the proxy --- unusually
likely to mislead.
The last is why the first two matter: a field that cannot yet tell whether an agent
succeeded is one whose every reported improvement is provisional, including the improvements
to how it measures.

{\footnotesize
\setlength{\bibsep}{2pt plus 1pt}
\bibliographystyle{abbrvnat}
\bibliography{../../refs}
}

\end{document}
